\section{Related Work}

\subsection{BLE Indoor Localization}

Indoor localization in healthcare and nursing environments has garnered significant attention, with BLE emerging as the industry standard due to cost-effectiveness, low power consumption, and unobtrusive deployment. Ayub et al.~\cite{ayub2025} compared six machine learning models—including Random Forest (RF), XGBoost, kNN, SVM, and MLP—on a BLE dataset across 10 rooms, reporting that RF achieves competitive F1-scores (0.91) at lower computational cost than deep alternatives. García-Paterna et al.~\cite{garcia2021} demonstrated room-level BLE localization using kNN with fingerprint averaging, achieving over 85\% accuracy in residential settings. Friedrich et al.~\cite{friedrich2024} proposed ensemble methods for BLE room detection in nursing facilities, reporting 2–3\% Macro F1 improvements.

However, the complex physical topology of nursing facilities introduces severe operational challenges. Factors such as human body shadowing, structural obstacles, and dynamic environmental changes cause multipath fading, resulting in highly erratic RSSI readings. Raw RSSI alone yields unacceptable room-level errors, necessitating advanced preprocessing and robust algorithmic interventions.

\subsection{RSSI Smoothing and Denoising}

Due to environmental interference, raw BLE RSSI signals are notoriously volatile. Common smoothing techniques include Simple Moving Average (SMA) and Kalman filters. However, SMA effectively reduces variance but inherently flattens signal peaks and blurs sudden RSSI shifts—features critical for identifying rapid spatial transitions, such as a caregiver walking through a doorway~\cite{hou2017}. Kalman filtering requires complex state-transition modeling that is difficult to tune for unpredictable human movement.

To overcome the loss of signal geometry, the Savitzky-Golay (SG) filter offers a compelling alternative. By applying local polynomial regression, SG smooths high-frequency noise while preserving the underlying shape, peaks, and boundaries of the original signal~\cite{savitzky1964}. Halder et al.~\cite{halder2015} applied SG as a baseline for ZigBee RSSI smoothing, demonstrating effective noise reduction while maintaining signal shape. Despite its proven efficacy in biosignal processing, SG filtering for BLE-based room-level \emph{classification} remains largely unexplored. This paper bridges this gap by introducing SG (window=5, polyorder=2) as a preprocessing step, achieving a 10.3\% improvement in Rare-F1 over raw RSSI.

\subsection{Class Imbalance in RSSI Data}

Data collected from real-world nursing facilities inherently exhibits extreme class imbalance. Caregivers spend disproportionately large amounts of time in communal areas (e.g., nurse station, 9,361 samples) compared to individual patient rooms (e.g., room 516, 23 samples). Conventional approaches to mitigate this often rely on synthetic data generation, most notably SMOTE~\cite{chawla2002}.

While SMOTE is effective on structured datasets, its application to highly stochastic RSSI data is problematic. Generating synthetic RSSI fingerprints introduces artificial noise and creates overlapping decision boundaries between physically distinct rooms, leading to overfitting. In our experiments, SMOTE produced the largest train-validation gap (0.1459) among all tested strategies. Threshold-based relabeling~\cite{garcia2024} has also been attempted, but cosine similarity proved ineffective when all rooms exhibit near-identical mean RSSI vectors ($\approx$\,$-$98\,dBm).

This study adopts a cost-sensitive approach via algorithmic class weighting (\texttt{balanced\_subsample}). By penalizing minority class misclassification within the classifier's loss function, the model accurately recognizes infrequent room visits while preserving the physical integrity of the original dataset. This configuration achieves the lowest overfit gap (0.0997) and a Weighted F1 of 0.8351, outperforming both SMOTE and relabeling.

\subsection{Temporal Features for Localization}

Traditional BLE fingerprinting relies on single-snapshot RSSI measurements, which suffer from spatial ambiguity—distant locations can exhibit identical RSSI signatures due to signal fluctuations. To address this, recent research has incorporated temporal features via sliding windows.~\citet{sensors2023} demonstrated that temporal fingerprint slices ($W=4$) from WiFi RSSI time series significantly improve accuracy over single-snapshot inputs.

While deep learning architectures like LSTM excel at sequential modeling, they incur high computational costs and require large training datasets, making them impractical for localized facility deployments. Furthermore, existing sliding window approaches often process raw or minimally filtered data, inheriting environmental noise.

The key novelty of this paper is the integration of SG filtering at the signal level with sliding window feature extraction ($k=5$, 4 statistics $\times$ 25 beacons $=$ 100 features) at the temporal level. This hybrid spatiotemporal approach empowers a lightweight Random Forest classifier to resolve spatial ambiguities without the overhead of deep learning models.
